
\documentclass{article}
\usepackage{colortbl}
\usepackage{makecell}
\usepackage{multirow}
\usepackage{supertabular}

\begin{document}

\newcounter{utterance}

\centering \large Interaction Transcript for game `cladder', experiment `full\_v1.5\_default', episode 7573 with Qwen3.5{-}4B.
\vspace{24pt}

{ \footnotesize  \setcounter{utterance}{1}
\setlength{\tabcolsep}{0pt}
\begin{supertabular}{c@{$\;$}|p{.15\linewidth}@{}p{.15\linewidth}p{.15\linewidth}p{.15\linewidth}p{.15\linewidth}p{.15\linewidth}}
    \# & \multicolumn{2}{c}{Player} && \multicolumn{2}{c}{Game Master} \\
    \hline

    \theutterance \stepcounter{utterance}  
    & & & \multicolumn{4}{p{0.6\linewidth}}{
        \cellcolor[rgb]{0.9,0.9,0.9}{
            \makecell[{{p{\linewidth}}}]{
                \texttt{\tiny{[P1$\langle$GM]}}
                \texttt{You are an expert in causal inference. The following question is not a typical commonsense query, but rather a meticulously designed question created by a professor specializing in causal inference, intended to assess the students' mastery of the course content.} \\
\texttt{Imagine a self{-}contained, hypothetical world with only the following conditions, and without any unmentioned factors or causal relationships: Pexu has a direct effect on hwax and rukz. Hwax has a direct effect on rukz. For those who are not pexu, the probability of rukz is 17\%. For those who are pexu, the probability of rukz is 80\%. For those who are pexu, would it be more likely to see rukz if the individual was not pexu?} \\
\texttt{Start your answer with 'Yes' or 'No', followed by additional reasoning or evidence to support your explanation.} \\
            }
        }
    }
     \\ \\

    \theutterance \stepcounter{utterance}  
    & \multicolumn{4}{p{0.6\linewidth}}{
        \cellcolor[rgb]{0.561,0.737,0.561}{
            \makecell[{{p{\linewidth}}}]{
                \texttt{\tiny{[P1$\rangle$GM]}}
                \texttt{No} \\
\\ 
\texttt{**Reasoning:**} \\
\\ 
\texttt{The question asks whether, for individuals who are **pexu**, the probability of observing **rukz** would be higher if those individuals were **not pexu**. This requires comparing two specific conditional probabilities:} \\
\\ 
\texttt{1.  **The actual scenario (Counterfactual for the subject):** The probability of rukz given that the individual is pexu.} \\
\texttt{    *   From the prompt: \$P(\textbackslash{}text\{rukz\} \textbackslash{}mid \textbackslash{}text\{pexu\}) = 80\textbackslash{}\%\$.} \\
\\ 
\texttt{2.  **The hypothetical scenario:** The probability of rukz given that the individual is **not** pexu.} \\
\texttt{    *   From the prompt: \$P(\textbackslash{}text\{rukz\} \textbackslash{}mid \textbackslash{}neg\textbackslash{}text\{pexu\}) = 17\textbackslash{}\%\$.} \\
\\ 
\texttt{The question essentially asks: Is \$P(\textbackslash{}text\{rukz\} \textbackslash{}mid \textbackslash{}neg\textbackslash{}text\{pexu\}) > P(\textbackslash{}text\{rukz\} \textbackslash{}mid \textbackslash{}text\{pexu\})\$?} \\
\texttt{Substituting the values: Is \$17\textbackslash{}\% > 80\textbackslash{}\%\$?} \\
\\ 
\texttt{The answer is clearly **No**. The probability of rukz is significantly lower (17\%) for those who are not pexu compared to those who are pexu (80\%). Therefore, if an individual is pexu, they are much more likely to see rukz than if they were not pexu. The presence of the intermediate variable **hwax** (which is affected by pexu and affects rukz) explains the causal mechanism behind this difference, but it does not alter the numerical comparison of the probabilities provided.} \\
            }
        }
    }
    & & \\ \\

    \theutterance \stepcounter{utterance}  
    & & & \multicolumn{2}{p{0.3\linewidth}}{
        \cellcolor[rgb]{0.9,0.9,0.9}{
            \makecell[{{p{\linewidth}}}]{
                \texttt{\tiny{[GM$|$GM]}}
                \texttt{No} \\
            }
        }
    }
    & & \\ \\

    \theutterance \stepcounter{utterance}  
    & & & \multicolumn{2}{p{0.3\linewidth}}{
        \cellcolor[rgb]{0.9,0.9,0.9}{
            \makecell[{{p{\linewidth}}}]{
                \texttt{\tiny{[GM$|$GM]}}
                \texttt{no} \\
            }
        }
    }
    & & \\ \\

    \theutterance \stepcounter{utterance}  
    & & & \multicolumn{2}{p{0.3\linewidth}}{
        \cellcolor[rgb]{0.9,0.9,0.9}{
            \makecell[{{p{\linewidth}}}]{
                \texttt{\tiny{[GM$|$GM]}}
                \texttt{game\_result = WIN} \\
            }
        }
    }
    & & \\ \\

\end{supertabular}
}

\end{document}
